\section{Гироскопы: классификация и физические принципы работы}
\label{sec:gyroscope}

Гироскоп — это устройство или прибор, способный измерять изменение углов ориентации связанного с ним тела либо самостоятельно определять и сохранять заданное направление в пространстве. В основе работы большинства гироскопов лежит фундаментальное свойство вращающегося тела — сохранение направления оси своего вращения при отсутствии воздействия внешних моментов сил \cite{babaeva1973giroskopy}.


Общая классификация гироскопов может быть проведена по различным признакам: по физическому принципу действия, функциональному назначению, уровню точности, конструктивному исполнению. Наиболее фундаментальной является классификация по физическому принципу действия, которая отражает эволюцию технологии:

\begin{itemize}
	\item Механические гироскопы (роторные).
	\item Оптические гироскопы (лазерные, волоконно-оптические).
	\item Вибрационные гироскопы (МЭМС, кварцевые).
\end{itemize}





